\sectioncentered*{Реферат}
\thispagestyle{empty}

%\MakeUppercase{\topicName}: дипломный проект/ \fio. -- Минск : БГУИР, \the\year{}, -- п.з. -- ~\pageref*{LastPage}~с., чертежей (плакатов) -- 6 л. формата А1.
\MakeUppercase{\topicName}: дипломный проект/ \fio. -- Минск : БГУИР, \the\year{}, -- п.з. -- ~\pageref*{LastPage}~с..

\vspace{4\parsep}

%Дипломный проект выполнен на 6 листах формата А1 с пояснительной запиской на~\pageref*{LastPage} страницах.
Пояснительная записка включает введение, аналитический раздел, раздел проектирования, раздел разработки, технико-экономическое обоснование, заключение и список использованных источников.

Ключевые слова: интеллектуальная система мониторинга, безопасность путешественников, управление рисками в поездках, обработка естественного языка, машинное обучение, реферирование текстов, машинный перевод, локализация данных, геоинформационная система, микросервисная архитектура, PostgreSQL, PostGIS, RabbitMQ.

Целью дипломного проекта является разработка интеллектуальной системы мониторинга рисков и обеспечения безопасности путешественников, предназначенной для автоматического сбора сведений о потенциально опасных событиях из открытых источников, семантического анализа новостных материалов, оценки географической зоны влияния угроз и оперативного информирования пользователей на выбранном языке.

Предметом исследования являются методы построения программных систем мониторинга туристических и командировочных рисков, включая подходы к обработке потоков событий, автоматическому определению языка источника, машинному переводу, реферированию текстов, классификации угроз, геопространственному анализу маршрутов и доставке локализованных уведомлений конечным пользователям.

Разработанный программный комплекс позволяет получать данные о событиях из внешних источников, выполнять первичную фильтрацию и очистку входного потока, извлекать текстовые описания инцидентов, определять язык исходного материала, переводить данные на поддерживаемые языки, формировать краткие аналитические сводки, рассчитывать зоны влияния событий и отображать их в пользовательском интерфейсе. Система также обеспечивает проверку маршрутов путешественников на пересечение с активными зонами риска и отправку уведомлений через веб-интерфейс, Push-канал и электронную почту.

В первом разделе пояснительной записки проведен анализ предметной области управления рисками в путешествиях, рассмотрены международные подходы к обеспечению принципа должной осмотрительности, исследованы существующие коммерческие аналоги, источники данных, архитектурные решения, методы определения языка, машинного перевода и реферирования текстов.

Во втором разделе выполнено проектирование функциональных требований, сценариев использования, микросервисной архитектуры, структуры базы данных, программных интерфейсов и конвейера интеллектуальной обработки событий. Особое внимание уделено хранению локализованных данных в JSONB-полях PostgreSQL, проектированию асинхронного взаимодействия сервисов и организации обработки текстов с учетом ограниченного набора языков, поддерживаемых моделями реферирования.

В третьем разделе описана программная реализация системы на языках Python и Java, приведены фрагменты кода модулей сбора данных, интеллектуального анализа, серверной логики, геопространственной обработки, многоканальных уведомлений и клиентского интерфейса. Представлены результаты разработки пользовательских экранов, карточек событий и email-уведомлений на русском и английском языках.

В четвертом разделе приведено технико-экономическое обоснование разрабатываемого программного продукта, включающее расчет затрат на создание системы, оценку цены программного продукта, экономического эффекта и показателей эффективности внедрения.

Результатом дипломного проектирования является программный комплекс интеллектуального мониторинга рисков и обеспечения безопасности путешественников, реализующий событийно-ориентированную микросервисную архитектуру, локальный контур обработки естественного языка, геопространственную модель оценки угроз и механизм доставки локализованных уведомлений пользователям.

\clearpage